\section{Instruksi MOV dan Aturan Operand}
\label{sec:mov}

Instruksi \asm{MOV} adalah instruksi transfer data yang paling dasar.
Instruksi ini menyalin nilai dari operand sumber ke operand tujuan.
Nilai sumber tidak berubah; yang berubah hanya tujuan. Dengan
demikian, \asm{MOV} lebih tepat dipahami sebagai ``menyalin'', bukan
``memindahkan''.

Bagian ini diringkas dari \cite{gunawan2019}, Bagian 5,
hlm.~35--44, dan \cite{sto2001}, BAB V, hlm.~21--24.

\subsection{Bentuk Penulisan}

Bentuk umum instruksi \asm{MOV} adalah:

\begin{center}
\lstasm|MOV tujuan, sumber|
\end{center}

Operand tujuan dan sumber dapat berupa register, lokasi memori, atau
nilai segera, dengan batasan tertentu. Bentuk yang diperbolehkan
dirangkum pada Tabel~\ref{tab:bentuk-mov}.

\begin{table}[H]
    \centering
    \caption{Bentuk instruksi MOV yang diperbolehkan.}
    \label{tab:bentuk-mov}
    \small
    \begin{tabular}{@{}p{3.0cm}p{4.2cm}p{3.0cm}p{4.0cm}@{}}
        \toprule
        \textbf{Bentuk} & \textbf{Contoh} &
        \textbf{Arah salinan} & \textbf{Keterangan} \\
        \midrule
        Register ke register & \lstasm|MOV AX, BX| &
        \reg{BX} $\rightarrow$ \reg{AX} &
        Ukuran kedua register harus sama \\
        \addlinespace
        Nilai segera ke register & \lstasm|MOV AL, 5| &
        5 $\rightarrow$ \reg{AL} &
        Nilai harus muat dalam ukuran register \\
        \addlinespace
        Memori ke register & \lstasm|MOV AX, nilai| &
        isi memori $\rightarrow$ \reg{AX} &
        Alamat diambil dari segmen data \\
        \addlinespace
        Register ke memori & \lstasm|MOV nilai, AX| &
        \reg{AX} $\rightarrow$ memori &
        Ukuran data harus sesuai \\
        \addlinespace
        Nilai segera ke memori & \lstasm|MOV nilai, 1234h| &
        nilai $\rightarrow$ memori &
        Ukuran data ditentukan direktif data \\
        \addlinespace
        Register ke segmen & \lstasm|MOV DS, AX| &
        \reg{AX} $\rightarrow$ \reg{DS} &
        Sumber harus register umum \\
        \bottomrule
    \end{tabular}
\end{table}

\subsection{Aturan Ukuran dan Jenis Operand}

Ada tiga aturan penting yang harus dipatuhi.

\begin{enumerate}
    \item \textbf{Ukuran operand harus sama.} Register 8 bit hanya
          dapat dipasangkan dengan data 8 bit, sedangkan register
          16 bit dengan data 16 bit. Instruksi
          \lstasm|MOV AL, 1234h| tidak sah karena \reg{AL} hanya
          menampung 8 bit.

    \item \textbf{Pemindahan memori ke memori tidak diperbolehkan.}
          Instruksi seperti \lstasm|MOV [BX], [SI]| harus dipecah
          menjadi dua langkah melalui register:
\begin{lstlisting}[language=MASM, numbers=none]
MOV AX, [SI]     ; memori -> register
MOV [BX], AX     ; register -> memori
\end{lstlisting}

    \item \textbf{Register segmen tidak dapat diisi nilai segera.}
          Nilai untuk \reg{DS}, \reg{SS}, atau \reg{ES} harus
          dipindahkan melalui register umum, seperti
          \lstasm|MOV AX, @DATA| lalu \lstasm|MOV DS, AX|.
\end{enumerate}

\subsection{Contoh Pemakaian}

\begin{lstlisting}[language=MASM, numbers=none]
MOV AX, 1234h       ; nilai segera ke register
MOV BX, AX          ; register ke register
MOV nilai, AX       ; register ke memori
MOV DX, nilai       ; memori ke register
MOV AL, 'A'         ; kode ASCII 'A' ke AL
MOV DS, AX          ; register umum ke register segmen
\end{lstlisting}

Instruksi \asm{MOV} tidak mengubah \eng{flag}. Karena itu, bila
program perlu memeriksa hasil operasi sebelumnya, \asm{MOV} aman
disisipkan di antara operasi dan pemeriksaan \eng{flag}.

\begin{peringatanbox}
\vspace{1.1em}
Dua kesalahan yang sering terjadi:
\begin{enumerate}
    \item \lstasm|MOV AL, 1234h| --- ukuran sumber melebihi ukuran
          tujuan.
    \item \lstasm|MOV [BX], [SI]| --- pemindahan memori ke memori
          tidak tersedia pada 8086.
\end{enumerate}
Keduanya harus diperbaiki dengan menyesuaikan ukuran atau memakai
register perantara.
\end{peringatanbox}

\begin{infobox}
\vspace{1.1em}
\asm{MOV} hanya menyalin data. Untuk menukar isi dua operand,
gunakan \asm{XCHG}; untuk membandingkan tanpa mengubah operand,
gunakan \asm{CMP}. Keduanya dibahas pada bagian berikut.
\end{infobox}
